\documentclass[10pt,a4paper]{amsart}
\usepackage[margin=19mm]{geometry}
\usepackage{amsmath,amssymb,mathtools}
\usepackage[scr=rsfs,cal=euler]{mathalfa}
\usepackage{xfrac}
\usepackage[unicode,hidelinks,bookmarks=false]{hyperref}
\newcommand{\ie}{i.e.\ }
\newcommand{\cf}{cf.\ }
\newcommand{\resp}{respectively\ }
\newcommand{\Subsection}{\subsection}
\providecommand{\nobreakdash}{\nobreak\hbox{-}}
\setlength{\emergencystretch}{2em}
\multlinegap0pt
\usepackage{bm}
\usepackage{bbm}
\usepackage{dsfont}
\usepackage{xspace}
\usepackage{cancel}
\usepackage{mathtools}
\usepackage{cleveref}
\usepackage{amsthm}
\makeatletter
\@ifundefined{theorem}{\newtheorem{theorem}{Theorem}}{}
\@ifundefined{lemma}{\newtheorem{lemma}[theorem]{Lemma}}{}
\@ifundefined{observation}{\newtheorem{observation}[theorem]{Observation}}{}
\makeatother
\def\C{\mathbb{C}}
\def\D{\mathbb{D}}
\def\T{\mathbb{T}}
\def\andspace{\quad\text{and}\quad}
\def\forallspace{\quad\text{for all}\quad}
\title[A sharp analysis of Root-MUSIC: locations of correct and ext]{A sharp analysis of Root-MUSIC: locations of correct and extraneous roots}
\author{Hana Huber, Weilin Li}
\date{Source: arXiv:2606.04003v2 (CC BY 4.0)}
\begin{document}
\maketitle

\noindent\emph{Source and licence.} A sharp analysis of Root-MUSIC: locations of correct and extraneous roots. Version arXiv:2606.04003v2 (CC BY 4.0), distributed under the Creative Commons Attribution 4.0 International licence, as recorded in the arXiv licence metadata for this exact version and shown on the abstract page https://arxiv.org/abs/2606.04003v2. Full rights notice: licenses/arxiv-2606.04003-CC-BY-4.0.txt.\par\smallskip

\noindent\textit{Editorial scope.} Counted unit: the Lemma whose source label is \texttt{lem:Ppowerdisk} in the article (source0.tex lines 1113--1119, source label \texttt{lem:Ppowerdisk}), delivered together with the article's own complete original proof of it (source0.tex lines 1121--1171, one \texttt{proof} environment of 51 lines). No mathematics is written, repaired or re-derived: statement and proof are the article's own text line for line. Required context retained uncounted: obs:3 (lines 599--610); lem:q'' (lines 1102--1107). Every definition, display and label of those lines is retained unchanged. Omitted from this package and not redistributed: 1--598, 611--1101, 1108--1112, 1120--1120, 1172--1696. Nothing inside a retained excerpt is omitted or altered: every retained body is delivered exactly as the article's lines are written, so no omission receipt of any kind accompanies this package. Citations: the retained excerpts carry no citation command at all, so no bibliography entry of the article is transcribed and no reference is redistributed. Reference adaptation: none was needed and none was made; every reference inside the delivered unit resolves inside it, so no unresolved reference or citation is produced. Printed slips kept literal: no printed word, symbol, hypothesis or number of the counted statement or of its proof was altered. Layout and class: the article's own class is \texttt{article} with packages including geometry, amsmath, appendix, amsthm, bbm, amssymb, setspace, color, float, dsfont, enumerate, hyperref, cleveref, graphicx; the wrapper is the lane's pinned \texttt{amsart} editorial wrapper at 10pt on A4, which restates the theorem-like environments the retained text needs and adds the article's own macro definitions that the retained text needs (\texttt{\textbackslash C}, \texttt{\textbackslash D}, \texttt{\textbackslash T}, \texttt{\textbackslash andspace}, \texttt{\textbackslash forallspace}), with extra wrapper packages (bm, bbm, dsfont, xspace, cancel, mathtools, cleveref). The delivered numbering is the unit's own. Assets: no retained span contains a figure environment, a graphics-inclusion command, a PDF-page inclusion, a table or a TikZ picture; no image file is copied into this package, the package declares no asset dependency and no image file is redistributed with it. Licence: version arXiv:2606.04003v2 of this article is distributed under the Creative Commons Attribution 4.0 International licence, as recorded in the arXiv licence metadata for this exact version and shown on the abstract page https://arxiv.org/abs/2606.04003v2. A full rights notice is delivered at licenses/arxiv-2606.04003-CC-BY-4.0.txt. Characters: every retained excerpt is pure ASCII, so the delivered engine, XeLaTeX with its default Latin Modern text font, reports no missing character.
\begin{observation}
	\label{obs:3}
	Suppose $m\geq s+1$. Then $Q$ has exactly $2s$ zeros on $\partial\D$, which are precisely $\{e^{i\theta_k}\}_{k=1}^s$, each with multiplicity two. Additionally $Q$ has at most $2m-2$ complex zeros, and they come in pairs with the following reflection symmetry: $w$ is a zero of $Q$ if and only if $1/\overline w$ is a zero. Finally, there is an algebraic polynomial $B$ of degree at most $m-1$ whose roots are all in $\overline \D$ such that
	\begin{equation}
		\label{eq:FR}
		Q(z)=B(z)\overline{B(1/\overline z)} \forallspace z\in \C.
	\end{equation}
    Consequently, we have the identity 
    $$
    Q(z)=\overline{Q(1/\overline z)} \forallspace z\in \C\setminus \{0\}.
    $$
\end{observation}

\begin{lemma} \label{lem:q''}
    Suppose $m\geq 100$ and $\Delta(\{\theta_k\}_{k=1}^s)\geq 8\pi/m$. There is an absolute constant $c_0\in (0,1/6)$ such that for all $k$, 
    $$
    q''(\theta_k) \geq c_0 m^2. 
    $$
\end{lemma}

\begin{lemma}\label{lem:Ppowerdisk}
    Suppose $m\geq 100$ and $\Delta(\{\theta_k\}_{k=1}^s)\geq 8\pi/m$. If $z\in D(w_k,c_0/(5m))$, then 
    $$
    |P(z)|
    \geq \frac 12 \left( 1 - \frac{8}{15} e^{2c_0/5}\right) c_0 m^2 |z-w_k|^2. 
    $$
\end{lemma}

\begin{proof}
    We plan to calculate the power series expansion of $P$ at its double root $w_k=e^{i\theta_k}$. 
    We have
	\begin{equation}
		\label{eq:Pseries}
		P(z)=\frac 12 P''(w_k)(z-w_k)^2 + \sum_{j=3}^{2m-2} \frac{P^{(j)}(w_k)}{j!} (z-w_k)^j. 
	\end{equation}
	By Bernstein's inequality for complex polynomials, we see that 
	$$
	\max_{z\in \partial\D} |P^{(j)}(z)|
	\leq (2m-2)^j \max_{z\in \partial\D} |P(z)|
	= (2m-2)^j \|q\|_{L^\infty(\T)}
	\leq (2m-2)^j. 
	$$ 
	Using this, we see that
	\begin{align*}
		\left| \sum_{j=3}^{2m-2} \frac{P^{(j)}(w_k)}{j!}(z-w_k)^j \right|
		&\leq 8(m-1)^3|z-w_k|^3 \, \sum_{j=3}^{2m-2} \frac{2^{j-3}(m-1)^{j-3} |z-w_k|^{j-3}}{j!} \\
		&\leq 8(m-1)^3|z-w_k|^3 \, \sum_{j=0}^\infty \frac{2^j(m-1)^j |z-w_k|^j}{(j+3)!} \\
        &\leq \frac{4}{3} (m-1)^3|z-w_k|^3 \, \sum_{j=0}^\infty \frac{2^j(m-1)^j |z-w_k|^j}{j!} \\
		&= \frac{4}{3} (m-1)^3 \, |z-w_k|^3 \, e^{2 (m-1)|z-w_k|}. 
	\end{align*}
    
    
    Additionally, \cref{obs:3} implies that $Q(w_k)=Q'(w_k)=0$ since each $w_k$ is a root with multiplicity two of $Q$. Then a calculation shows that 
    $$
    P''(w_k)=w_k^{m-1}Q''(w_k) \andspace q''(\theta_k)=-w_k^2 \, Q''(w_k).
    $$
    By \cref{lem:q''}, there is a $c_0\in (0,1/6)$ such that $q''(\theta_k)\geq c_0m^2$. Thus,
	$$
	|P''(w_k)|
	=|Q''(w_k)|
	=|q''(\theta_k)|
	\geq c_0 m^2.
	$$
	Thus, inserting the above inequality into \eqref{eq:Pseries}, we have
	\begin{align*}
		|P(z)|
		&\geq \frac 12 |P''(w_k)||z-w_k|^2 - \left| \sum_{j=3}^{2m-2} \frac{P^{(j)}(w_k)}{j!}(z-w_k)^j \right| \\
		&\geq \frac 12 c_0 m^2 |z-w_k|^2- \frac 4 3(m-1)^3|z-w_k|^3 e^{2(m-1)|z-w_k|} \\
		&\geq \frac 12 m^2 |z-w_k|^2 \left( c_0 - \frac 83 m|z-w_k| e^{2 m|z-w_k|}\right).
	\end{align*}
    By assumption, $|z-w_k|\leq c_0/(5m)$, and so 
    $$
    c_0-\frac 8 3m|z-w_k| e^{2 m|z-w_k|}
    \geq c_0-\frac{8}{15} e^{2c_0/5} c_0.
    % \leq \frac{4}{15} e^{1/15} c
    % < \frac 3 {10} c.  
    $$
    Combining the above completes the proof. 
\end{proof}

\end{document}
